\documentclass[a4paper,11pt]{article}
        \usepackage[top=15mm,bottom=18mm,left=16mm,right=16mm]{geometry}
        \usepackage{fontspec}
        \usepackage{xeCJK}
        \setmainfont{Times New Roman}
        \setCJKmainfont{Yu Gothic}
        \setmonofont{Consolas}
        \setCJKmonofont{Yu Gothic}
        \usepackage{booktabs}
        \usepackage{longtable}
        \usepackage{tabularx}
        \usepackage{array}
        \usepackage{enumitem}
        \usepackage{hyperref}
        \usepackage{listings}
        \usepackage{xcolor}
        \hypersetup{
          colorlinks=true,
          linkcolor=blue,
          urlcolor=blue,
          pdftitle={live\_wifi ROS2 launch 入口},
          pdfauthor={Codex}
        }
        \lstset{
          basicstyle=\ttfamily\small,
          breaklines=true,
          columns=fullflexible,
          keepspaces=true,
          frame=single,
          framerule=0.2pt,
          rulecolor=\color{black},
          showstringspaces=false
        }
        \setlength{\parskip}{0.42em}
        \setlength{\parindent}{1em}
        \renewcommand{\arraystretch}{1.15}
        \title{03. live\_wifi ROS2 launch 入口}
        \author{solar\_ws0129-main}
        \date{2026-07-29}
        \begin{document}
        \maketitle

        \section*{対象}
        \begin{tabularx}{\linewidth}{>{\raggedright\arraybackslash}p{0.18\linewidth} X}
        \toprule
        項目 & 内容 \\
        \midrule
        ファイル & \path|launch/solar_race_live_wifi.launch.py| \\
        種別 & ROS 2 launch Python \\
        区分 & launch \\
        役割 & WiFi テレメトリ、風補正、live MPC 運用をまとめて起動する launch 入口。 \\
        起動文脈 & live\_wifi モード起動時に ros2 launch される。 \\
        \bottomrule
        \end{tabularx}

        \section*{このファイルを読む理由}
        WiFi テレメトリ、風補正、live MPC 運用をまとめて起動する launch 入口。

        \begin{itemize}[leftmargin=1.6em]
        \item profile\_yaml を launch 引数として受ける。
\item 共通 live node 群に WiFi bridge と wind correction を追加する。
\item live forecast の raw/corrected CSV の入口でもある。
        \end{itemize}

        \section*{呼び出し関係}
        \subsection*{どこから呼ばれるか}
        \begin{itemize}[leftmargin=1.6em]
        \item \path|scripts/solar_control.sh|
\item \path|SolarSim.ps1|
        \end{itemize}

        \subsection*{次に読むと理解が進むファイル}
        \begin{itemize}[leftmargin=1.6em]
        \item \path|mpc_solarcar/live_launch.py|
\item \path|mpc_solarcar/telemetry_text_bridge_node.py|
\item \path|mpc_solarcar/wind_correction_node.py|
        \end{itemize}

        \subsection*{同一リポジトリ内の主要依存}
        \begin{itemize}[leftmargin=1.6em]
        \item \path|mpc_solarcar/live_launch.py|
\item \path|mpc_solarcar/solar_profile.py|
        \end{itemize}

        \section*{ソース冒頭の把握}
        モジュール docstring または先頭意図の要約:

        モジュール docstring は明示されていない。

        \subsection*{代表コード断片}
        \begin{lstlisting}[language=Python]
nodes = build_live_nodes(profile_path, cfg, use_wifi=True)
    if bool(wifi_cfg.get("enabled", True)):
        nodes.append(
            Node(
                package="mpc_solarcar",
                executable="telemetry_text_bridge_node",
                name="telemetry_text_bridge_node",
                parameters=[
                    {
                        **{key: value for key, value in wifi_cfg.items() if key != "enabled"},
                        "solar_power_gain_to_pack": float(
                            model_cfg.get("solar_measurement_gain_to_pack", 1.0)
                        ),
                    }
\end{lstlisting}


        \section*{主要構造の読み取り}
        主要関数は generate\_launch\_description。 launch から起動する Node action は 2 件。

        \subsection*{主要クラス・関数}
            \begin{longtable}{p{0.07\linewidth} p{0.16\linewidth} p{0.25\linewidth} p{0.40\linewidth}}
            \toprule
            行 & 種別 & 名前 & 補足 \\
            \midrule
            12 & function & \path|_setup| & args: context \\
62 & function & \path|generate_launch_description| &  \\
            \bottomrule
            \end{longtable}

        \subsection*{ファイルを上から読んだときの定義順}
            \begin{longtable}{p{0.07\linewidth} p{0.84\linewidth}}
            \toprule
            行 & 内容 \\
            \midrule
            12 & 関数 \_setup を定義する。 \\
62 & 関数 generate\_launch\_description を定義する。 \\
            \bottomrule
            \end{longtable}

        \subsection*{import 群の詳細}
            \begin{longtable}{p{0.07\linewidth} p{0.33\linewidth} p{0.50\linewidth}}
            \toprule
            行 & import 文 & このファイルで読む理由 \\
            \midrule
            1 & \path|from __future__ import annotations| & 型ヒントの遅延評価を有効にし、自己参照型や forward reference を安全に扱うため。 このファイル内での使用位置は少ないか、間接利用である。 \\
3 & \path|from launch import LaunchDescription| & launch が実行すべき action 群をまとめるため。 このファイル内での主な使用位置は L63。 \\
4 & \path|from launch.actions import DeclareLaunchArgument, OpaqueFunction| & DeclareLaunchArgument や OpaqueFunction など launch action を使うため。 このファイル内での主な使用位置は L65, L66。 \\
5 & \path|from launch.substitutions import LaunchConfiguration| & launch 引数の実行時値を参照するため。 このファイル内での主な使用位置は L13。 \\
6 & \path|from launch_ros.actions import Node| & ROS 2 ノード起動 action を launch から記述するため。 このファイル内での主な使用位置は L25, L41。 \\
8 & \path|from mpc_solarcar.live_launch import build_live_nodes, live_forecast_paths| & live系 node 構成ビルダ から build\_live\_nodes, live\_forecast\_paths を読み込み、このファイルの内部処理を分担させるため。 実体ファイルは mpc\_solarcar/live\_launch.py。 このファイル内での主な使用位置は L20, L22。 \\
9 & \path|from mpc_solarcar.solar_profile import get_section, load_profile| & profile YAML 読込と検証 から get\_section, load\_profile を読み込み、このファイルの内部処理を分担させるため。 実体ファイルは mpc\_solarcar/solar\_profile.py。 このファイル内での主な使用位置は L14, L15, L16, L17, L18。 \\
            \bottomrule
            \end{longtable}

        \section*{関数・クラスを上から順に解説}
        \subsection*{L12 関数 \path|_setup|}
            \begin{itemize}[leftmargin=1.6em]
              \item 定義: \path|_setup(context)|
              \item docstring: 明示されていない。
              \item このブロックが直接呼ぶ主な関数・メソッド: LaunchConfiguration, Node, append, bool, build\_live\_nodes, float, get, get\_section, items, live\_forecast\_paths, load\_profile, perform
              \item 戻り値の要点: nodes
            \end{itemize}
            \begin{enumerate}[leftmargin=1.8em]
            \item profile\_yaml に LaunchConfiguration('profile\_yaml').perform(context) の結果を代入する。
\item (profile\_path, cfg) に load\_profile(profile\_yaml) の結果を代入する。
\item live\_cfg に get\_section(cfg, 'live') の結果を代入する。
\item wifi\_cfg に get\_section(live\_cfg, 'wifi\_bridge') の結果を代入する。
\item wind\_cfg に get\_section(live\_cfg, 'wind\_model') の結果を代入する。
\item model\_cfg に get\_section(cfg, 'model') の結果を代入する。
\item (\_base\_csv, raw\_forecast\_csv, corrected\_csv) に live\_forecast\_paths(profile\_path, cfg) の結果を代入する。
\item nodes に build\_live\_nodes(profile\_path, cfg, use\_wifi=True) の結果を代入する。
\item 条件 bool(wifi\_cfg.get('enabled', True)) を判定し、真なら内部処理を行う。
\item 条件 bool(wind\_cfg.get('enabled', True)) を判定し、真なら内部処理を行う。
            \end{enumerate}

\subsection*{L62 関数 \path|generate_launch_description|}
\begin{itemize}[leftmargin=1.6em]
  \item 定義: \path|generate_launch_description()|
  \item docstring: 明示されていない。
  \item このブロックが直接呼ぶ主な関数・メソッド: DeclareLaunchArgument, LaunchDescription, OpaqueFunction
  \item 戻り値の要点: LaunchDescription([DeclareLaunchArgument('profile\_yaml', default\_value='config/solar/bwsc\_2027\_demo.yaml'), OpaqueFunction(function=\_setup)])
\end{itemize}
\begin{enumerate}[leftmargin=1.8em]
\item LaunchDescription([DeclareLaunchArgument('profile\_yaml', default\_value='config/solar/bwsc\_2027\_demo.yaml'), OpaqueFunction(function=\_setup)]) を返す。
\end{enumerate}





        \subsection*{launch から起動するノード}
            \begin{longtable}{p{0.07\linewidth} p{0.30\linewidth} p{0.50\linewidth}}
            \toprule
            行 & executable & package / name \\
            \midrule
            25 & \path|telemetry_text_bridge_node| & package=mpc\_solarcar, name=telemetry\_text\_bridge\_node \\
41 & \path|wind_correction_node| & package=mpc\_solarcar, name=wind\_correction\_node \\
            \bottomrule
            \end{longtable}





        \section*{処理の流れ}
        \begin{enumerate}[leftmargin=1.7em]
        \item launch 引数を受け取り、profile を読み込む。
\item Node action を動的に組み立てる。
\item LaunchDescription として ROS 2 launch へ返す。
        \end{enumerate}

        \section*{このファイルの位置づけ}
        この資料群では、\path|launch/solar_race_live_wifi.launch.py| を
        \textbf{launch}の中核として扱う。
        したがって、ここで扱う処理は単独で完結せず、
        前段の profile / launch / telemetry と後段の planner / logger / report へ繋がる。

        \end{document}
